\chapter*{Resumen}
\addcontentsline{toc}{chapter}{Resumen}

Esta investigación analiza las trayectorias académicas de estudiantes de las licenciaturas en
Matemáticas y Física de la Facultad de Ciencias (UNAM) con el objetivo de caracterizar,
explicar y predecir el rezago académico mediante métodos estadísticos multivariados y de
aprendizaje automático. A partir del \emph{Vector Excalibur} —un conjunto de indicadores
semestrales de regularidad gregoriana, índice de esfuerzo, promedio natural y promedio
real— se construyó una arquitectura analítica de tres enfoques paralelos: exploratorio
(PCA), inferencial (Análisis de Supervivencia) y predictivo (Random Forest y Regresión
Regularizada).

El análisis de componentes principales redujo el espacio de 32 variables a cinco componentes
que explican el $88.5\,\%$ y $90.9\,\%$ de la varianza en Matemáticas y Física,
respectivamente, revelando una gradación continua entre trayectorias eficientes y rezagadas
que se manifiesta desde el segundo semestre. El análisis de supervivencia (Kaplan-Meier,
Log-Rank y Cox) confirmó una diferencia estadísticamente significativa entre carreras
($\chi^2 = 5.85$, $p = 0.016$) e identificó la \textbf{paradoja del perfeccionismo}: los
estudiantes con alta regularidad y promedio modesto egresan sistemáticamente antes que los
de alta calificación pero baja regularidad, con un \emph{Hazard Ratio} de la regularidad de
$\text{HR} = 12.70$ ($p < 0.005$) frente a un efecto del promedio no significativo
($p = 0.08$).

Los modelos de clasificación Random Forest —validados mediante validación cruzada estratificada
y bootstrapping ($B = 2{,}000$)— alcanzaron una exactitud del $65.5\,\%$
(IC~95\,\%: $[64.1\,\%, 66.9\,\%]$) con información de los primeros cuatro semestres y del
$78.2\,\%$ (IC~95\,\%: $[77.0\,\%, 79.4\,\%]$) con la trayectoria completa; la diferencia
es estadísticamente significativa (McNemar: $\chi^2 = 827.70$,
$p = 5.13 \times 10^{-182}$). La regresión regularizada Lasso sobre la submuestra sin censura
($N = 3{,}157$; generaciones $\leq 2016$) explicó el $75.5\,\%$ de la varianza del semestre
de egreso ($R^2 = 0.755$) con un Error Absoluto Medio de $0.957$ semestres, siendo la
regularidad en el octavo semestre el predictor dominante en todos los modelos ($\beta = -2.445$,
$t = -13.64$, $p < 10^{-40}$).

Los resultados sugieren que la \textbf{inercia de inscripción continua} —y no la excelencia
en la calificación— es el mecanismo estadístico dominante del egreso oportuno, con implicaciones
directas para el diseño de sistemas de alerta estudiantil basados en regularidad en lugar de
promedio.

\bigskip
\noindent\textbf{Palabras clave:} rezago académico, análisis de supervivencia, random forest,
regresión regularizada, paradoja del perfeccionismo, inercia académica, alerta temprana.

% ──────────────────────────────────────────────────────────────────────────────
\chapter*{Abstract}
\addcontentsline{toc}{chapter}{Abstract}

This thesis analyzes the academic trajectories of undergraduate students in Mathematics and
Physics at the Faculty of Sciences (UNAM) with the goal of characterizing, explaining, and
predicting academic delay using multivariate statistical methods and machine learning. Drawing
on the \emph{Vector Excalibur} —a set of semestral indicators covering enrollment regularity,
effort index, natural GPA, and real GPA— a three-pronged analytical architecture was designed:
exploratory (PCA), inferential (Survival Analysis), and predictive (Random Forest and
Regularized Regression).

Principal Component Analysis reduced the 32-variable space to five components explaining
$88.5\,\%$ and $90.9\,\%$ of variance in Mathematics and Physics, respectively, revealing
a continuous gradient between on-time and delayed trajectories visible as early as the second
semester. Survival analysis (Kaplan-Meier, Log-Rank, and Cox) confirmed a statistically
significant difference between programs ($\chi^2 = 5.85$, $p = 0.016$) and identified the
\textbf{Perfectionist Paradox}: students with high enrollment regularity but modest grades
graduate systematically earlier than those with high grades but low regularity, with a
regularity Hazard Ratio of $\text{HR} = 12.70$ ($p < 0.005$) versus a non-significant
effect of GPA ($p = 0.08$).

Random Forest classifiers —validated via stratified cross-validation and bootstrapping
($B = 2{,}000$)— reached $65.5\,\%$ accuracy (95\,\% CI: $[64.1\,\%, 66.9\,\%]$) using
only the first four semesters, and $78.2\,\%$ (95\,\% CI: $[77.0\,\%, 79.4\,\%]$) with the
full eight-semester trajectory; the gain is statistically significant (McNemar:
$\chi^2 = 827.70$, $p = 5.13 \times 10^{-182}$). Lasso regression on the uncensored subsample
($N = 3{,}157$; cohorts $\leq 2016$) explained $75.5\,\%$ of graduation-semester variance
($R^2 = 0.755$) with a Mean Absolute Error of $0.957$ semesters, with eighth-semester
enrollment regularity as the dominant predictor across all models ($\beta = -2.445$,
$t = -13.64$, $p < 10^{-40}$).

The findings suggest that \textbf{continuous enrollment inertia} —not academic excellence
as measured by grades— is the statistically dominant mechanism of on-time graduation, with
direct implications for student early-warning systems that should prioritize regularity over
GPA as the primary risk signal.

\bigskip
\noindent\textbf{Keywords:} academic delay, survival analysis, random forest, regularized
regression, perfectionist paradox, academic inertia, early warning system.
